\chapter{Cas d'égalité des triangles}

\begin{itemize}
\it Deux triangles $ABC$ et $A'B'C'$ ont $AB=A'B'$, $\widehat{B}=\widehat{B'}$ et $\widehat{C}=\widehat{C'}$. On mène les hauteurs $(AH)$ et $(A'H')$.\begin{enumerate}
\item Comparer les triangles rectangles $ABH$ et $A'B'H'$. Conséquences ?
\item Comparer les triangles $ACH$ et $A'C'H'$, puis les triangles $ABC$ et $A'B'C'$. 
\item Énoncer le cas d'égalité non classique correspondant.
\end{enumerate}
\it Comparer deux triangles isocèles dans chacun des cas suivants : 
\begin{enumerate}
\item Les bases sont égales et les angles au sommet égaux.
\item Les hauteurs relatives à la base sont égales ainsi que les angles à la base. 
\item Énoncer les cas d'égalité correspondants. 
\end{enumerate}
\it Dans un triangle $ABC$, on prolonge la médiane $(AM)$ d'une longueur $MD=AM$. 
\begin{enumerate}
\item Comparer les triangles $BMD$ et $CMA$, puis évaluer les côtés du triangle $ABD$ par rapport aux côtés $[AB]$, $[AC]$ et à la médiane $(AM)$. 
\item Comparer deux triangles $ABC$ et $A'B'C'$ ayant deux côtés respectivement égaux ainsi que la médiane relative au troisième. Énoncer le cas d'égalité correspondant. 
\end{enumerate}
\it Deux triangles $ABC$ et $A'B'C'$ ont un côté égal $BC=B'C'$ ainsi que les hauteurs $(AH)$ et $(A'H')$ et les médianes $(AM)$ et $(A'M')$. 
\begin{enumerate}
\item Comparer les triangles $AMH$ et $A'M'H'$.
\item On superpose ces deux triangles. En déduire l'égalité des triangles $ABC$ et $A'B'C'$ et énoncer le cas d'égalité correspondant. 
\end{enumerate}
\it Démontrer que dans un triangle isocèle $ABC$ de base $[BC]$. \begin{enumerate}
\item Les médianes relatives aux côtés égaux sont égales.
\item Les bissectrices intérieures des angles $\widehat{B}$ et $\widehat{C}$ sont égales. 
\item Les hauteurs $(BH)$ et $(CK)$ sont égales. Étudier la réciproque. 
\end{enumerate}
\it On prend sur les côtés d'un angle $\widehat{xAy}$ deux points $B$ et $C$ ($AB\neq AC)$. 
La bissectrice de $\widehat{xAy}$ et la médiatrice de $[BC]$ se coupent en $D$. \begin{enumerate}
\item Comparer $DB$ et $CD$ puis les distances $DE$ et $DF$ de $D$ aux côtés de $\widehat{xAy}$.
\item Comparer les triangles $DBE$ et $DCF$, puis les angles $\widehat{BDC}$ et $\widehat{EDF}$.
\item Montrer que $BE$ et $FC$ sont égaux à la demi-différence de $AB$ et $AC$. 
\end{enumerate}
\it Soit un triangle isocèle $ABC$, dans lequel la base $[BC]$ est inférieure aux côtés égaux $[AB]$ et $[AC]$. 
On prolonge $[AB]$ et $[BC]$ de longueurs $BD$ et $CE$ égales à la différence $AB-BC$. \begin{enumerate}
\item Montrer que $BE=AC$. Puis comparer les triangles $ACE$ et $EBD$.
\item Montrer que $\widehat{ADE}=\frac12\left(\widehat{AED}+\widehat{BAC}\right)$.
\end{enumerate}
\it Soient deux points $A$ et $B$ équidistants d'une même droite $(xy)$. On désigne par $M$ et $N$ les pieds des perpendiculaires menées de $A$ et $B$ sur $(xy)$ et par $O$ le milieu de $[MN]$. \begin{enumerate}
\item Comparer les triangles $OAM$ et $OBN$. Conséquences ? 
\item On suppose que $A$ et $B$ soient de part et d'autre de $(xy)$. Montrer que $O$ est aussi le milieu de $[AB]$.
\item On suppose $A$ et $B$ du même côté de $(xy)$. Montrer que la médiatrice de $[AB]$ est la médiatrice de $[MN]$.
\end{enumerate}
\it On considère un triangle isocèle $ABC$ dans lequel 
la médiatrice du côté $[AC]$ coupe le prolongement de la base $[BC]$ en un point $D$. On prolonge $[DA]$ d'une longueur $AE=BD$.\begin{enumerate}
\item Montrer que le triangle $DAC$ est isocèle. Conséquences ?
\item Comparer les triangles $ABD$ et $CAD$. Que peut-on dire du triangle $CDE$ ?
\end{enumerate}
\it La médiatrice du côté $[AB]$ du triangle isocèle $ABC$ coupe la base $[BC]$ (ou son prolongement) en $D$. Le cercle de centre $B$ passant par $D$ recoupe $(AD)$ en $E$. 
\begin{enumerate}
\item Montrer que les triangles $DAB$ et $BDE$ sont isocèles. Comparer les longueurs $AD$ et $BE$ puis les angles $\widehat{ACD}$ et $\widehat{BAE}$ ainsi que les angles $\widehat{ADC}$ et $\widehat{BEA}$.
\item En utilisant le résultat de l'exercice \textbf{25} démontrer l'égalité des triangles $ACD$ et $BEA$ puis que $AE=CD$.
\end{enumerate}
\it \begin{enumerate}
\item Deux triangles $ABC$ et $A'B'C'$ ont $AB=A'B'$, $\widehat{B}+\widehat{B'}= 180^o$ et $\widehat{C}= \widehat{C'}$. Montrer que les hauteurs $AH$ et $A'H'$ sont égales, puis que $AC=A'C'$. 
\item Réciproquement si $AB=A'B'$, $AC=A'C'$ et $\widehat{B}+\widehat{B'}= 180^o$, les angles $\widehat{C}$ et $\widehat{C'}$ sont égaux. (On pourra amener $A'B'C'$ sur $ABC$ de telle sorte que $[A'B']$ coïncide avec $[AB]$ et que $C'$ vienne sur $(CB)$.)
\end{enumerate}
\it Soit un triangle isocèle $ABC$ de base $[BC]$. On
construit une demi-droite $[Bx)$ intérieure à l'angle $\widehat{ABC}$ et une demi-droite $[Cy)$ intérieure à l'angle $\widehat{ACB}$ de telle sorte que les angles $\widehat{ABx}$ et $\widehat{ACy}$ soient égaux. $(Bx)$ et $(Cy)$ se coupent en $M$.
\begin{enumerate}
\item Soient $H$ et $K$ les pieds des perpendiculaires menées de $A$ à $(Bx)$ et à $(Cy)$. Comparer les triangles $ABH$ et $ACK$, puis les segments $[AH]$ et $[AK]$.
\item Établir que la droite $(AM)$ est bissectrice extérieure du triangle $BMC$.
\end{enumerate}
\it Les bissectrices intérieurs des angles en $B$ et $C$ du triangle $ABC$ se coupent en $I$, tandis que leurs bissectrices extérieures se coupent en $J$.\begin{enumerate}
\item Montrer que $I$ et $J$ sont équidistants des trois droites $(AB)$, $(AC)$ et $(BC)$.
\item Démontrer que les trois points $A$, $I$ et $J$ sont alignés.
\end{enumerate}
\it Dans un quadrilatère convexe $ABCD$, on a : $AB=AD$ et $\widehat{B}+\widehat{D}=180^o$. On prolonge $[CB]$ d'une longueur $BE=CD$. \begin{enumerate}
\item Comparer les triangles $ABE$ et $ADC$. Nature du triangle $ACE$. 
\item Démontrer que $(CA)$ est bissectrice intérieure de l'angle en $C$ du quadrilatère.
\end{enumerate}
\it On considère un triangle $ABC$ dans lequel $AB>AC$. La médiatrice de $[BC]$ coupe en $M$ la bissectrice intérieure de l'angle en $A$. Le point $M$ se projette perpendiculairement en $H$ sur $[AB]$ et en $K$ sur le prolongement de $[AC]$.
\begin{enumerate}
\item Comparer les triangles rectangles $MBH$ et $MCK$, puis les triangles rectangles $AMH$ et $AMK$. Démontrer que $AH=AK$ et $BH=CK$. 
\item En déduire que : $AH = \frac12(AB+AC)$ et $BH = \frac12(AB-AC)$.
\end{enumerate}
\it Dans le triangle $ABC$ tel que $AB<AC$, la médiatrice de $[BC]$ coupe en $D$ le côté $[AC]$, en $I$ la bissectrice intérieure de l'angle en $A$ et en $J$ sa bissectrice extérieure.
\begin{enumerate}
\item Montrer que le triangle $DBC$ est isocèle. Que 
représente la droite $(IJ)$ pour l'angle en $D$ du triangle $ADB$ ?
\item Montrer que chacun des points $I$ et $J$ est équidistant des trois droites $(AB)$, $(AD)$ et $(BD)$ puis que $(BI)$ et $(BJ)$ sont bissectrices intérieure
et extérieure de l'angle $\widehat{ABD}$. 
\item Démontrer que les angles $\widehat{ABJ}$ et
$\widehat{ACJ}$ sont égaux tandis que les angles 
$\widehat{ABI}$ et $\widehat{ACI}$ sont supplémentaires.
\end{enumerate}
\end{itemize}








